% FAC-EN-U0008
% French source: fac_body.tex lines 269--273.
% Authority: PDF physical p. 10; printed p. 205 (no. 9 complete).

\medskip
\phantomsection\label{fac:I:1:no:9}
\textbf{9. Direct sum of two sheaves.}

Let $\mathcal{A}$ be a sheaf of rings, and let $\mathcal{F}$ and $\mathcal{G}$ be two sheaves of $\mathcal{A}$-modules. For every $x\in X$, form the module $\mathcal{F}_x+\mathcal{G}_x$, the \textit{direct sum} of $\mathcal{F}_x$ and $\mathcal{G}_x$; an element of $\mathcal{F}_x+\mathcal{G}_x$ is a pair $(f,g)$, with $f\in\mathcal{F}_x$ and $g\in\mathcal{G}_x$. Let $\mathcal{H}$ be the disjoint union of the $\mathcal{F}_x+\mathcal{G}_x$ as $x$ ranges over $X$. We may identify $\mathcal{H}$ with the subset of $\mathcal{F}\times\mathcal{G}$ formed by the pairs $(f,g)$ such that $\pi(f)=\pi(g)$. If $\mathcal{H}$ is endowed with the topology induced by that of $\mathcal{F}\times\mathcal{G}$, it is immediately verified that $\mathcal{H}$ is a sheaf of $\mathcal{A}$-modules. It is called the \textit{direct sum} of $\mathcal{F}$ and $\mathcal{G}$ and is denoted by $\mathcal{F}+\mathcal{G}$. Every section of $\mathcal{F}+\mathcal{G}$ over an open set $U\subset X$ has the form $x\mapsto(s(x),t(x))$, where $s$ and $t$ are sections of $\mathcal{F}$ and $\mathcal{G}$ over $U$. In other words, $\Gamma(U,\mathcal{F}+\mathcal{G})$ is isomorphic to the direct sum $\Gamma(U,\mathcal{F})+\Gamma(U,\mathcal{G})$.

The definition of the direct sum extends inductively to a finite number of sheaves of $\mathcal{A}$-modules. In particular, the direct-sum sheaf of $p$ sheaves isomorphic to a given sheaf $\mathcal{F}$ will be denoted by $\mathcal{F}^p$.
